\section{Recovering the arithmetic statement behind step 17}
\label{sec:step17}

\paragraph{Source contract and provenance.}
Alp\"oge's \emph{Remark}, p.~4, step~17 and footnote~6, names
$H^1(\mathcal O_E[1/\ell],\mathbb Q_\ell/\mathbb Z_\ell(-e))$
and a weakening of Vandiver, but supplies neither quantifiers nor a theorem
number for that weakening. The attached map identifies candidate modules only.
We recovered those modules from the public mirror
\url{https://github.com/CocoSgt/riemann-hypothesis-lean}, at commit
\texttt{7cbf7a0c32fca1badf05931879a37c3f8a2b03b4}.
The retrieved \texttt{src/Main.lean} is byte-identical to the supplied file;
its SHA-256 is recorded in the evidence manifest. Git blob identifiers were
checked for the manifest-based downloads. These are source inspections, not
a new Lean build or a claim to have authenticated the entire mirror. Under
the prompt's stipulated theorem contract, we use the recovered unconditional
Alp\"oge theorem below. The arithmetic translation and consequences are the
new deductions checked here.

In this section a source path beginning with \texttt{AV/} abbreviates
\path{src/Development/AutomorphicArithmetic/AutomorphicVarieties/}.
The source facts have precise declaration locators; no theorem is inferred
merely from a file name or from a proposition's definition.

\subsection{The local conditions and the fields}

Put $D_\ell(e)=\mathbb Q_\ell/\mathbb Z_\ell(-e)$ and, for a number field $E$,
\[
 \mathcal S_\ell(E,e)=
 \ker\left(H^1_{\mathrm{et}}(\operatorname{Spec}\mathcal O_E[1/\ell],D_\ell(e))
       \longrightarrow\prod_{v\mid\ell}H^1(E_v,D_\ell(e))\right).
 \tag{S1}
\]
Thus the selected classes are unramified outside $\ell$ and locally trivial
at every place above $\ell$. They are a subgroup of the displayed $H^1$,
not that entire group. The source definitions are
\path{src/Formalism/GaloisCohomology/Statement.lean}, declarations
\texttt{LData}, \texttt{qzTw}, \texttt{Sha}, and \texttt{shaLen}.

For a finite extension coefficient ring $O\subset\overline{\mathbb Q}_\ell$
which is finite over $\mathbb Z_\ell$ and is a DVR, write
\[
 D_O=O\otimes_{\mathbb Z_\ell}(\mathbb Q_\ell/\mathbb Z_\ell),
 \qquad e(O)=\operatorname{length}_O(O/\ell O).
 \tag{S2}
\]
For a Dirichlet character $\psi$ modulo $N$ with values in $O^\times$, define
$\mathcal V_{\ell,N}(e,\psi)$ to be the subgroup of
$H^1(G_{\mathbb Q},D_O(\psi^{-1}\chi_\ell^{-e}))$
whose restriction is zero on every finite inertia group and on every
decomposition group at a prime dividing $\ell N$. Here $\chi_\ell$ is the
cyclotomic character. Set
\[
 v_{\ell,N}(e,\psi)
 =\operatorname{length}_O\mathcal V_{\ell,N}(e,\psi)/e(O).
 \tag{S3}
\]
In (S3) the quotient of lengths is used only when the module has finite
length. Literally, the source's \texttt{Vn} applies
\texttt{ENat.toNat} to its extended-natural lengths (which sends
infinity to zero); every estimate used below simultaneously proves
finiteness, so (S3) then gives exactly its value. These are
\texttt{Coef}, \texttt{Vlen}, and \texttt{Vn} in
\path{AV/HeightAndVanishing/DirichletPrimeToLPartMulLPowerPart.lean}, together
with \texttt{QO}, \texttt{eO}, \texttt{shaFam}, and \texttt{ShaU} in
\path{AV/GaloisLibrary/SelmerTransport/ConjStableShaFam.lean}.
The scalar character is defined in
\path{AV/CharacterLevelBounds/ShaFamMono.lean}, \texttt{rho}.

\begin{proposition}[K17-FIELD-DESC: the field interpretation]
Let $e>0$, let $\ell$ be odd, and suppose $E/\mathbb Q$ is a cyclic
extension of degree prime to $\ell$, cut out by $\psi$. Suppose $O$ contains
the values of $\psi$. Restriction gives an isomorphism
\[
 \mathcal V_{\ell,N}(e,\psi)
 \simeq
 \left(O\otimes_{\mathbb Z_\ell}\mathcal S_\ell(E,e)\right)^{\psi},
 \tag{S4}
\]
where the right side is the $\psi$-eigenspace for the natural action of
$\operatorname{Gal}(E/\mathbb Q)$ on cohomology, including its action on
the Tate coefficients. The same description applies to character components
inside a larger abelian extension of degree prime to $\ell$.
\end{proposition}
\begin{proof}
Write $\Delta=\operatorname{Gal}(E/\mathbb Q)$. All positive-degree
cohomology of $\Delta$ with $\ell$-primary coefficients vanishes: it is
annihilated by $|\Delta|$, which acts invertibly on those coefficients.
Inflation--restriction consequently identifies the global $H^1$ on the
left with the indicated eigenspace of $H^1(G_E,D_O(-e))$. The same argument
applies to decomposition and inertia groups, since their quotient groups
have order dividing $|\Delta|$. Hence restriction detects the vanishing of
each local restriction.

An unramified local class at a place $w\nmid\ell$ of $E$ lies in
\[
 H^1_{\mathrm{unr}}(E_w,D_O(-e))
 =D_O/(\operatorname{N}w^{-e}-1)D_O=0.
\]
The last equality holds because $\operatorname{N}w^{-e}-1$ is a nonzero
element of $\mathbb Q_\ell$ and multiplication by any nonzero scalar from
$\mathbb Z_\ell$ is surjective on the divisible module $D_O$.
Thus all extra decomposition conditions away from $\ell$, including
those dividing $N$, are automatic after the inertia conditions. The local
conditions at $\ell$ are exactly (S1). Restriction therefore gives (S4).
Finite flatness of $O$ over $\mathbb Z_\ell$ commutes with the cochain
description and the eigenspace idempotent, so no scalar-extension defect
remains. This proves the assertion.
\end{proof}

For CM fields the source's ``plus'' notation is twist-dependent. If $c_E$
is complex conjugation, then
\[
 \mathcal S_\ell(E,e)^+
 =\{x\in\mathcal S_\ell(E,e):c_E x=(-1)^e x\}.
 \tag{S5}
\]
It is not always the $+1$ eigenspace. See
\path{AV/LengthEqAddEigen.lean}, \texttt{ShaPlus}, \texttt{ShaMinus},
\texttt{shaLen\_eq\_add}, and
\path{AV/PartsSign.lean}, \texttt{mk\_mem\_ShaPlus\_iff}.
The direct sum of these two eigenspaces uses that $2$ is invertible;
the source treats $\ell=2$ separately.

The fields used in the larger twisting argument are not arbitrary number
fields. The source's \texttt{Host} fixes a Dirichlet character $\chi$ of
modulus $m$, an integer $a\ge2$, and primes $q\equiv3\pmod8$,
$d\equiv7\pmod{16}$ with $(qd,m)=1$. They determine
\[
 E=F_\chi\,\mathbb Q(\mu_{2^{a+2}})^+
      \mathbb Q(\sqrt{-q})\,\mathbb Q(\sqrt{-d}),
 \qquad N=\operatorname{lcm}(2^{a+2},m)qd,
 \tag{S6}
\]
where $F_\chi$ is the fixed field of $\ker\chi$. This is the fixed-field
description encoded by \texttt{IsCompField} and
\texttt{isCompField\_ofChar} in \path{AV/StmtFields.lean}, and by
\texttt{Host.H} in
\path{AV/HeightAndVanishing/DirichletPrimeToLPartMulLPowerPart.lean}.
Indeed the four congruence conditions in that definition fix respectively
the real $2$-power cyclotomic field, $F_\chi$, and the two quadratic fields;
the intersection of the corresponding subgroups fixes their compositum.
The selected character components have $\psi(-1)=(-1)^e$.

For completeness, the exact finite-parameter organization is in
\path{AV/HeightAndVanishing/LayerLe.lean}, where \texttt{PointHV}
is the proposition defined to assert the following statement:
for each $\chi$ and function $j:\mathbb N\to\mathbb N$, there are
$C,r_1,r_2\in\mathbb N$, with $r_1,r_2$ prime, such that, for every host
with character $\chi$, both auxiliary primes smaller than $q,d$, every
$t\ge2$ satisfying \texttt{Host.Big}, and every $B\ge0$ satisfying
the specified \texttt{LayerLe} bounds at the finite set \texttt{Bchi},
the weighted sum of the lengths in (S5), over all odd primes and
$1\le e<t$, is bounded by $(Nt)^C+CB$.
Here \texttt{Host.Big} means $2t+1<q,d$ and each of
$m,\operatorname{ord}\chi,p_K(\chi),h_K(\chi),p_1(\chi)$ is less than
$q,d$. The finite exceptional set is
\[
 B_\chi=\{3\}\cup
 \{\text{odd prime divisors of }
   (\operatorname{ord}\chi)m p_K(\chi)h_K(\chi)p_1(\chi)\}.
\]
The three auxiliary quantities are the least primes/class number specified
in \path{AV/StmtFields.lean}, \texttt{pK}, \texttt{hStar}, \texttt{pOne}.
We record this definition to identify the fields and classes in the
note. The closed source proof composition in the following host-family
subsection supplies its theorem inputs; the stated layer bounds remain
explicit premises. The deductions about $\mathbb Z$ below use the unconditional,
quantified theorem next, and do not assume an unverified numerical value
for any of these constants or layer bounds.

\subsection{The unconditional character theorem}

The exact admissibility condition \texttt{GD} is in
\path{AV/HeightAndVanishing/Counts.lean}. In ordinary notation it says:
$e\ge1$, $N$ is even, $\ell$ is odd, the order of $\psi$ is prime to
$\ell$, the factor $\psi_\ell$ supported at the $\ell$-part of the
conductor satisfies $\psi_\ell^2=1$, and
$\psi(-1)=(-1)^e$. There is an auxiliary imaginary quadratic field $K$
such that
\[
 \ell\nmid |D_K|h_K,\quad (|D_K|,N)=1,\quad
 K\not\subset\mathbb Q(\mu_\ell)F_{\psi/\psi_\ell},\quad
 (\ell-1)\nmid2e. \tag{S7}
\]
The third condition is exactly the source's subgroup noncontainment
\texttt{eta}; it has not been silently deleted. The source definition
of \texttt{chiP} agrees with $\psi/\psi_\ell$ under its prime-to-$\ell$
order hypothesis. In particular these conditions force $\ell\ge5$.

\begin{sourcefact}[Alp\"oge: K17-S17, the recovered quantitative theorem]
For every imaginary quadratic $K$ there exist $c_5,A\in\mathbb N$ with
the following property. Fix $e,N$, a subgroup
$H\subset(\mathbb Z/N\mathbb Z)^\times$, and distinct primes $p,p'$
not dividing $30N|D_K|$. For each prime $\ell$ in a finite set $L$ choose
a permitted coefficient ring $O_\ell$. If at least one character with
kernel $H$ satisfies (S7) and the preceding admissibility conditions at
each $\ell\in L$, then all the associated lengths are finite and
\[
 \sum_{\ell\in L}\log\ell
   \sum_{\ker\psi=H}v_{\ell,N}(e,\psi)
 \le 2c_5\varphi([(\mathbb Z/N\mathbb Z)^\times:H])
          (N(e+1))^A(\log p+\log p'). \tag{S8}
\]
The character sum is over the characters valued in $O_\ell$; existence
of one such faithful character provides its whole orbit.

Exact locators: \path{AV/UnitaryThreeGroups/KatoUse.lean},
\texttt{CorGK}; and \path{AV/SmallClassNumbers/CorGKOf.lean},
\texttt{UpperTheory.V.SmallClassNumbers.corGK}, an unconditional theorem.
\end{sourcefact}

\paragraph{Proof reconstruction and the cited inputs.}
The arithmetic input \texttt{PropGK} supplies, for an auxiliary prime $p$,
a nonzero integer $n_p$ such that
\[
 \log|n_p|\le c_5\varphi([G:H])(N(e+1))^A\log p,
 \qquad
 \sum_{\ker\psi=H}v_{\ell,N}(e,\psi)\le2v_\ell(n_p)
 \quad(\ell\ne p). \tag{S9}
\]
Its source proof is \texttt{propGK} in
\path{AV/UnitaryThreeGroups/OrbEquiv.lean}. It combines
\texttt{OneChar} (the Selmer-length estimate by rank-three/rank-two
automorphic congruences) with \texttt{CongruenceNumberRankThree}
(a nonzero norm of Frobenius-trace differences). The latter provides the
integer and its archimedean size, while the former bounds its local
valuation. The faithful-character orbit has size $\varphi([G:H])$.
Their unconditional inputs are assembled in
\texttt{SmallClassNumbers.corGK\_of}, \texttt{litN}, and \texttt{corGK}
of the file cited above: local and global Galois cohomology, the Kato
input, adjoint vanishing, automorphy, and congruence-number estimates.
These are cited Alp\"oge theorems under the source contract; they are
not independent new automorphy proofs in this paper.

Here is the complete elementary passage from (S9) to (S8).
Use (S9) twice, with $p$ and $p'$. For each $\ell$ at least one applies,
so the nonnegative length sum is at most
$2v_\ell(n_p)+2v_\ell(n_{p'})$. For a nonzero integer $n$,
\[
 \sum_{\ell\in L}v_\ell(n)\log\ell\le\log|n|
\]
because the product of the indicated prime powers divides $|n|$.
Summing and using the two archimedean estimates proves (S8), including
the auxiliary-prime cases. This is precisely
\path{AV/UnitaryThreeGroups/CorGSumLog.lean}, \texttt{corGK}.
The complete source proofs of this passage and of the orbit/norm
translation were read in this run.

\subsection{The rational-field specialization}

\begin{theorem}[K17-S17-Q]
There exist absolute $C,A\in\mathbb N$ such that for every positive even
integer $e$ and every finite set of primes $L$ satisfying
$\ell\ge5$ and $(\ell-1)\nmid2e$,
\[
 \mathcal S_\ell(\mathbb Q,e)\text{ is finite},\qquad
 \sum_{\ell\in L}\log\ell\,
   \operatorname{length}_{\mathbb Z_\ell}\mathcal S_\ell(\mathbb Q,e)
 \le C(e+1)^A. \tag{S10}
\]
In particular the group is zero whenever $\log\ell>C(e+1)^A$
and the displayed prime conditions hold.
\end{theorem}
\begin{proof}
In (S8) take $K=\mathbb Q(\sqrt{-3})$, $N=2$, $H=G$, the trivial
character, and auxiliary primes $p=7,p'=11$. The discriminant is $-3$
and the class number is one. The character conditions are immediate;
evenness of $e$ gives the parity condition. For $\ell\ge5$ the field
$K$ is not contained in $\mathbb Q(\mu_\ell)$, since the latter is
unramified at $3$ whereas $K$ is ramified there. This checks (S7).
All these checks are also made by \texttt{gd} in
\path{AV/SmallClassNumbers/ClassNumberThree.lean}.

Choose the bottom coefficient ring. It is isomorphic to
$\mathbb Z_\ell$ by \texttt{exists\_coef} in
\path{AV/SmallClassNumbers/LenEqZeroIff.lean}. Then $D_O=D_\ell(0)$
and $e(O)=1$, since $O/\ell O\simeq\mathbb F_\ell$ is a simple
$O$-module. The extra local condition at $2$ is automatic by the
unramified local calculation in the proof of K17-FIELD-DESC.
Thus the length in (S8) is the length in (S10).

The right side of (S8) is
$2c_5 2^A(e+1)^A(\log7+\log11)$, bounded by
$36c_5 2^A(e+1)^A$. Take $C=36c_5 2^A$.
This is the proof \texttt{bound\_of} in
\texttt{ClassNumberThree.lean}; its unconditional input and conclusion
are \texttt{corGK} and \texttt{bound} in
\texttt{SmallClassNumbers/CorGKOf.lean}.
For the last assertion use $L=\{\ell\}$. A nonzero finite
$\mathbb Z_\ell$-module has length at least one, contradicting
the strict inequality on $\log\ell$. The source independently states
this cocycle formulation as \texttt{Large} and proves
\texttt{large\_of} in \texttt{LenEqZeroIff.lean}.
\end{proof}

\paragraph{What is and is not numerical here.}
The source has also a proposition \texttt{CorG} with a displayed
$e^7N^{18}|D_K|^5$ shape. A definition of that proposition is not a
proof of it. We have used the established \texttt{CorGK}, with its
existential exponent, throughout. Neither ``100ish'' in footnote~6
nor an unspecified constant is a certified degree range. The next
section gives every vanishing and cyclicity conclusion that actually
follows from (S10), and identifies the remaining numerical gate.
